\section{The cubic multiplication at the character boundary}
\label{app:cubic-witness}

The proof of Theorem~\ref{thm:explicit-stable-character-polynomials}
leaves one possible multiplicity-one cubic source. We exhibit a
highest vector whose image survives the quadratic quotient. The
dual-number specialization below makes all unwanted moments vanish
and leaves only one term of that image.

\begin{lemma}[A nonzero cubic moment multiplication]
\label{lem:cubic-witness}
For $a\ge3$, put $n=2a$, $D=a+3$, and
\[
 \tau=(a-1,1),\qquad\mu=(a,a-1,1),\qquad\lambda=(a,2,1).
\]
For $\dim V\ge3$, the cubic multiplication map
\[
 \operatorname{Sym}^3V\otimes
       \operatorname{Hom}_{S_n}(S^\mu,Q_{n,k}[a])
      \longrightarrow\operatorname{Hom}_{S_n}(S^\mu,Q_{n,k}[a+3])
\]
is nonzero, and hence injective, on its multiplicity-one
$\mathbb S_\lambda V$ summand. Here $Q_{n,k}$ denotes the quotient
by linear and quadratic moments.
\end{lemma}
\begin{proof}
It is enough to work in three rows $x,y,z$, because the map is a
natural transformation of polynomial functors. Set $m=a-1$ and
choose disjoint label blocks of sizes $3,2,\ldots,2,1$, with $a-2$
blocks of size two. Define
\[
 F=\det\begin{pmatrix}1&x_0&y_0\\1&x_1&y_1\\1&x_2&y_2\end{pmatrix}
       \prod_{j=1}^{a-2}(x_{2j+1}-x_{2j+2}).
\]
The unused label is $2a-1$. This polynomial has row highest weight
$(m,1)$, since $\sum_i x_i\partial_{y_i}$ kills its determinant.
It belongs to label $\mu$. To justify this last assertion directly,
its column antisymmetry forces every label constituent to be
 dominated by $\mu$, hence to have first row at most $a$; degree $a$
forces first row at least $a$. At first row exactly $a$, the
row multiplicity in degree $a$ is $\mathbb S_\nu V$ for label
$(a,\nu)$, by the fully occupied, linear part of Schur--Weyl and
Pieri branching (equivalently, Theorem~\ref{thm:stable-augmented-multiplicities}
in its linear degree). The row highest weight $(a-1,1)$ therefore selects
exactly $\mu$. All invariant differential operators annihilate this
first occurrence, so $F$ survives the quadratic quotient.

Write $p_{ijk}=\sum_s x_s^iy_s^jz_s^k$, and put
\[
 A=\sum_s y_s\partial_{x_s},\qquad
 B=\sum_s z_s\partial_{x_s},\qquad
 C=\sum_s z_s\partial_{y_s}.
\]
The row weight $(a,2,1)$ and the two simple raising equations
determine the following combination of cubic moments and lowered
copies of $F$:

\begin{align*}
 G={}&p_{300}(AB-A^2C)F-(m-2)p_{210}BF
 +(2m-1)p_{210}ACF\\
 &-(m+1)p_{201}AF-(m^2-1)p_{120}CF
 +(m^2-1)p_{111}F.
\end{align*}
It has weight $(a,2,1)$ and is highest. For completeness, if
$X=\sum_sx_s\partial_{y_s}$ and
$Z=\sum_sy_s\partial_{z_s}$, then
\[
 XAF=(m-1)F,\quad XBF=-CF,\quad XACF=mCF,
\]
\[
 X(AB-A^2C)F=(m-2)BF-(2m-1)ACF,
\]
while $ZBF=ZACF=AF$, $ZCF=F$, $ZAF=0$, and
$Z(AB-A^2C)F=0$. The corresponding identities on the cubic moments
show $XG=ZG=0$. These are the simple raising operators of
$\operatorname{GL}_3$. Every operator used preserves the label
$\mu$, and Pieri's rule gives the indicated row summand once.

We show that $G$ survives the quadratic quotient by an exact
specialization over $K(\zeta)[\epsilon]/(\epsilon^2)$, where $\zeta$
is a primitive $2a$th root of unity. The frequencies $1,a-1,a$
sum to $2a$ for $xyz$. Every linear or quadratic frequency is
nonzero modulo $2a$, except that of $z^2$, whose coefficient
vanishes because $\epsilon^2=0$.
Set, for $0\le i<2a$,
\[
 x_i=\zeta^i,\qquad y_i=\zeta^{(a-1)i},\qquad
 z_i=\epsilon\zeta^{ai}.
\]
Every linear or quadratic moment is zero. For quadratic moments
containing two $z$'s this follows from $\epsilon^2=0$; every other
vanishing is a geometric sum with exponent not divisible by $2a$.
The same observation gives
\[
 p_{300}=p_{210}=p_{201}=p_{120}=0,
                      \qquad p_{111}=2a\epsilon.
\]
All factors of $F$ are nonzero. Its first determinant is
$\zeta^2-\zeta^{-2}$, which is nonzero for $a\ge3$, and each
remaining factor is a difference of distinct roots of unity.
Consequently
\[
                           G=2a(m^2-1)\epsilon F\ne0.
\]
The specialization kills the quadratic ideal, so $G$ is nonzero
in $Q_{n,3}$. Characteristic-zero scalar descent and functoriality
prove the assertion for all $V$.
\end{proof}

